\documentclass[12pt]{article}
%% BUGS:
%% - The \begin{problem} ... \end{problem} is only used in exams??

\usepackage{url, graphicx}
\PassOptionsToPackage{normalem}{ulem} % magic
\usepackage{href-ul}

% page layout
\usepackage[letterpaper, textheight=9.5in, textwidth=5.5in]{geometry}
\usepackage{setspace}
\setstretch{1.08}
\pagestyle{plain}
\sloppy\sloppypar\raggedbottom\frenchspacing\thispagestyle{empty}

% problem set header
\newcommand{\psname}{Problem Set}
\newcommand{\startps}[1]{\section*{Computational Physics --- \psname~{#1}}}

% problem formatting
\usepackage{amsthm}
\theoremstyle{definition}
\newcommand{\problemname}{Problem}
\newtheorem{problem}{\problemname}
\newcommand{\probref}[1]{\problemname~\ref{#1}}
\theoremstyle{remark}
\newcommand{\notename}{Note}
\newtheorem{note}{\notename}[problem]
\newcommand{\noteref}[1]{\textit{\notename}~\ref{#1}}
\newcounter{subpart}[problem]
\newcommand{\subpart}{\par\addvspace{0.75ex}\refstepcounter{subpart}\textsl{(\alph{subpart})}~}

% words
\newcommand{\foreign}[1]{\textsl{#1}}
\newcommand{\vs}{\foreign{vs}}

% code
\newcommand{\code}[1]{\texttt{\detokenize{#1}}}
\usepackage{xcolor}
% Define custom colors for code highlighting
\definecolor{codegreen}{rgb}{0,0.6,0}
\definecolor{codegray}{rgb}{0.5,0.5,0.5}
\definecolor{codepurple}{rgb}{0.58,0,0.82}
\definecolor{backcolour}{rgb}{0.98,0.98,0.99}
\usepackage{listings}
\lstdefinestyle{pythonstyle}{
    backgroundcolor=\color{backcolour},   
    commentstyle=\color{codegreen},
    keywordstyle=\color{magenta},
    numberstyle=\tiny\color{codegray},
    stringstyle=\color{codepurple},
    tabsize=4
}
\lstset{
    basicstyle=\small\ttfamily,
    columns=fixed,
    breaklines=false,
    keepspaces=true,
    showspaces=false,
    showstringspaces=false,
    showtabs=false,
    showlines=*,
    style=pythonstyle,
    language=Python,
    literate={~}{{\url{~}}}1
}

% math
\usepackage{amsmath}
\newcommand{\dd}{\mathrm{d}}
\newcommand{\e}{\mathrm{e}}

% primary units
\newcommand{\rad}{\mathrm{rad}}
\newcommand{\kg}{\mathrm{kg}}
\newcommand{\m}{\mathrm{m}}
\newcommand{\s}{\mathrm{s}}

% secondary units
\renewcommand{\deg}{\mathrm{deg}}
\newcommand{\km}{\mathrm{km}}
\newcommand{\cm}{\mathrm{cm}}
\newcommand{\mm}{\mathrm{mm}}
\newcommand{\ft}{\mathrm{ft}}
\newcommand{\mi}{\mathrm{mi}}
\newcommand{\AU}{\mathrm{AU}}
\newcommand{\ns}{\mathrm{ns}}
\newcommand{\h}{\mathrm{h}}
\newcommand{\yr}{\mathrm{yr}}
\newcommand{\N}{\mathrm{N}}
\newcommand{\J}{\mathrm{J}}
\newcommand{\eV}{\mathrm{eV}}
\newcommand{\W}{\mathrm{W}}
\newcommand{\Pa}{\mathrm{Pa}}

% derived units
\newcommand{\mps}{\m\,\s^{-1}}
\newcommand{\mph}{\mi\,\h^{-1}}
\newcommand{\mpss}{\m\,\s^{-2}}

% random units
\newcommand{\au}{\mathrm{a.u.}}

% draft comments
\newcommand{\hogg}[1]{{\color{blue}\textbf{#1}}}

\begin{document}

\startps{5}
Due Friday 2026 October 9 at 18:00 on Brightspace.
Be sure to explicitly give credit where it is due, on all problems.
Also credit buddies, web pages, and LLMs (as per course policy).
Failure to give explicit credits will result in points deducted.

\begin{problem}
  \subpart
  Write a very short Python program \code{make_hello()} that checks (using functions in the \code{os} package) whether the file \code{hello.txt} exists.
  If it exists, do nothing.
  If it doesn't exist, create it by opening a text file, writing the word \code{hello} into it, and closing it.
  \subpart
  Write a very short Python program \code{make_goodbye()} that checks whether both files \code{hello.txt} and \code{goodbye.txt} exist.
  First: If \code{hello.txt} doensn't exist, run \code{make_hello()} and then re-run \code{make_goodbye()}.
  That's recursive.
  Then: If \code{goodbye.txt} doesn't exist, create it and put into it the contents of \code{hello.txt} plus the additional word \code{goodbye}.
  \subpart
  Run \code{make_hello()} and \code{make_goodbye()} in order and check that the contents of \code{goodbye.txt} is what you expect.
  \subpart
  Now modify \code{make_goodbye()} so that it checks that the file modification date of \code{goodbye.txt} is later than the modification data of \code{hello.txt}.
  If it isn't, have your code delete \code{goodbye.txt} and then re-run \code{make_goodbye()}.
  That's also recursive.
  Test that this functionality works as expected, perhaps by running \code{make_goodbye()} and then \code{make_hello()} and then \code{make_goodbye()} again.
  \subpart
  Why did Prof Hogg make this \problemname?
\end{problem}

\begin{note}
  Learn how to use Python \code{with} and use it for your file creation and file reading operations.
\end{note}

\begin{note}
  When your code makes a file, have it print ``I'm making file [filename]!''.
  This will make it way easier to tell what your code is choosing to do!
  It will also be way easier for Iraj to figure out what your code did.
\end{note}

\begin{problem}
  In exoplanet astronomy, we solve the following equation over and over again:
  \begin{align}
    M &= E - e\,\sin E ~,
  \end{align}
  where $M$ is the ``mean anomaly'' (in radians),
  $E$ is the ``eccentric anomaly'' (in radians),
  and $e$ is the eccentricity (dimensionless) of an orbit.
  \subpart
  Write a function \code{M_from_E(E, e)} that computes and returns $M$ (note that this is very simple!)
  \subpart
  Write a function \code{E_from_M(M, e)} that computes and returns $E$ (note that this requires a root-finding).
  You can use binary search or else Newton's method or even \emph{relaxation}.
  Your code should, internally, call \code{M_from_E} as part of the root-finding operation; that makes sure that these two codes are inverses of each other.
  \subpart
  Make a scatter plot showing $M$ for a uniform grid of 12 angles $E$ in $0\leq E 2\pi$ for $e=0., 0.5, 0.75, 0.95$.
  Give each set of points a label and a legend on the figure.
  Make sure you label all your axes correctly.
  \subpart
  Make another scatter plot showing $E$ for a uniform grid of 12 angles $M$.
  Do your two plots agree in the relevant ways?
\end{problem}

\begin{note}
  If you use Newton's method, you might also want to write a function \code{dM_dE(E, e)} that outputs the derivative $\left.\dd M/\dd E\right|_{E,e}$.
\end{note}

\begin{note}
  If you want to learn some Python Jax, you could write \code{dM_dE} by writing \code{M_from_E} in Jax and then applying the \code{jax.grad} function.
\end{note}

\begin{note}
  The two functions you write should be inverses of each other.
  How would you write a sensible, non-trivial test of that?
  If you want to do more here, write that test and see if your functions pass it.
  What other tests should the functions pass?
\end{note}

\begin{problem}
  Newman, \textit{Computational Physics, 2ed}, Exercise~6.18.
\end{problem}

\begin{problem}
  Newman, \textit{Computational Physics, 2ed}, Exercise~6.21.
\end{problem}

\begin{note}
  You don't have to do part (e) of this problem, but if you love it and want to, we will admire you for doing it.
\end{note}

\begin{note}
  Don't get too distracted by any zero eigenvalues you find.
  Interpret them the way you did in \textsl{Physics~III}.
  The ``lowest frequency'' is related to the lowest non-zero eigenvalue.
\end{note}

\end{document}
